\subsubsection*{冲量把短暂作用与动量变化联系起来}

牛顿第二定律写成 $\dd\vec p/\dd t=\vec F$ 后，在 $t_1$ 到 $t_2$ 上积分得
\begin{equation}
 \Delta\vec p=\vec p_2-\vec p_1
 =\int_{t_1}^{t_2}\vec F\,\dd t\equiv\vec J.
\end{equation}
冲量是力对时间的累积，方向由矢量积分决定；在 $F$--$t$ 图上，每个分量的冲量等于相应的
有向面积。碰撞力可能很大且随时间快速变化，但只要碰撞前后的动量可测，便无需知道力的每个
瞬时值。平均力 $\bar{\vec F}=\vec J/\Delta t$ 只保证总冲量相同，并不表示真实力恒定。

对质点系逐个应用动量定理并相加，成对内力因第三定律消去：
\begin{equation}
 \vec P_2-\vec P_1=\int_{t_1}^{t_2}\vec F_{\rm ext}\,\dd t.
\end{equation}
碰撞、爆炸等短过程里，重力等有限外力的冲量常比接触冲量小得多，可以在这一时间尺度上
忽略；这不是说外力不存在，而是做了量级比较。若某方向外冲量不可忽略，只有其余方向的动量
分量可能守恒。把过程时间无限放大或把很小的外力累积很久，原近似就会失效。

连续介质问题可用“单位时间通过控制面的动量”理解。质量流率为 $\dot m=\rho Av$ 的稳定射流，
若速度由 $\vec v_1$ 变为 $\vec v_2$，维持该变化所需合力为
\begin{equation}
 \vec F=\dot m(\vec v_2-\vec v_1).
\end{equation}
管壁所受力方向与此相反。弯管、喷嘴、雨滴落入车厢和水柱冲击挡板看似不同，本质都是在选定
时间内统计流入、流出和留在系统内的动量。画清控制面和速度方向，比死记某个压力公式可靠。
